\section{Approach}
\label{sec:approach}

\subsection{Route 1: causal discovery over SAE features}

\begin{enumerate}[nosep]
  \item Run the forecaster over held-out trajectories and record intermediate
        (processor) activations at every step.
  \item Train an SAE on those activations to obtain a sparse feature basis.
  \item For each trajectory, read out each feature's activation at every step,
        giving multivariate time series $\{a_f(t)\}_{f=1}^{F}$.
  \item Run PCMCI+ over the feature series (TSCI and DYNOTEARS serve as
        cross-checks) to recover a lagged directed graph among features.
  \item Score the graph. On SAVAR this means precision, recall, F1, sign accuracy,
        and threshold-free AUROC/AUPRC against the known graph. On GraphCast, where
        no ground truth exists, it means consistency with known physical
        relationships and with interventional edits in the style
        of~\cite{macmillan2025}.
\end{enumerate}

Agreement among the three discovery methods on SAVAR, where we can score them,
calibrates how much to trust each of them later on GraphCast, where we cannot.

\subsection{Route 2: parameter decomposition}

The second route applies VPD (Sec.~\ref{sec:pd-background}) to the forecaster's
shared message-passing MLPs. Each weight matrix becomes a set of rank-one
components with a per-node gate; averaging the gate over inputs and reshaping to
the grid gives every component a spatial usage map. Since the components themselves
are shared across all nodes, exactly as in GraphCast's processor, any position
specificity in those maps had to emerge through the gate.

The two routes answer different questions and inherit different weaknesses. SAE
features live in activation space and come with a time axis, so they support
claims about dynamics, but their causal role must be established after the fact
and the time axis brings aliasing problems (below). VPD components live in weight
space, their importance is an ablation by construction, and no time axis is
involved; what they cannot directly show is how the represented quantities evolve.

\subsection{Do SAE features support temporal causal claims?}
\label{sec:temporal}

An SAE is trained per activation vector, so a feature is an atemporal descriptor
of a single state, and it is fair to ask whether causal discovery, which is about
time, is even well posed on such features. It is, because the time axis never
comes from the SAE. It comes from the trajectory: the SAE supplies the variables,
the forecaster's evolution across steps supplies the dynamics, and reading a
feature's activation across steps gives an ordinary time series.

The argument comes with caveats, which we treat as real limitations and return to
in Sec.~\ref{sec:open}. A feature has to keep its meaning from one step to the
next, or its time series is incoherent. On GraphCast the collection regime
matters: free-running rollouts accumulate error and drift off-distribution, so the
series is non-stationary and conditional-independence tests are confounded,
whereas teacher-forced trajectories (consecutive single steps on observed states)
stay approximately stationary. We therefore collect on teacher-forced trajectories
unless closed-loop behaviour is itself under study. Discovery on internal features
also recovers the causal structure of the \emph{model's} computation, which need
not coincide with the atmosphere's; for the question ``what did the model learn?''
the model's structure is the right target, but the claim should be stated as such.
And the 6-hour step bounds the resolvable lag: anything faster aliases into
contemporaneous dependence. Section~\ref{sec:subsample} measures this effect and
what PCMCI+ recovers of it.
